%% toc: P1 8 8.4
\section{Proofs for the motive letters instantiated by the colon-equals sign}
\label{sec:annex-p01-w02c-p01-018-motives}

\begin{annexitem}{annex:p01-w02c-p01-018-motive-readings}{thm:p01-w02c-p01-018-motive-readings}{The readings of a motive letter with one value and with several values}
\begin{proof}
The lines \(o_1,\ldots,o_n\) are all the lines of \(R\) with \(p\) on the
left of \(:=\), and they occur in this order.  A free occurrence of \(p\)
is one that is not bound by a statement; by hypothesis every free
occurrence lies in some application \(\alpha_i\), and since the
applications are pairwise disjoint it lies in exactly one of them.  Each
application \(\alpha_i\) begins at its line \(o_i\)
(\cref{def:p01-w02c-p01-018-motive}, item~2), so every free occurrence in
\(\alpha_i\) stands at or after \(o_i\).

\emph{Item~1.}  Read every \(o_i\) in role~(b).  By item~1 of
\cref{def:p01-w02c-p01-018-motive-reading}, the line \(o_i\) gives \(p\)
the value \(\lambda\vec z.\,s_i\) for the application \(\alpha_i\), and a
free occurrence of \(p\) inside \(\alpha_i\) means that value.  No line of
role~(a) has \(p\) on its left, because the only lines with \(p\) on the
left are the \(o_i\), all read in role~(b); so no region-wide name \(p\)
exists, and the value of \(o_i\) is the only meaning that the reading gives
to an occurrence in \(\alpha_i\).  Since each free occurrence lies in
exactly one \(\alpha_i\), it has exactly one meaning, namely
\(\lambda\vec z.\,s_i\).

\emph{Item~2.}  Let \(n=1\) and read \(o_1\) in role~(a).  By item~2 of
\cref{def:p01-w02c-p01-018-motive-reading}, \(p\) is a name for
\(\lambda\vec z.\,s_1\) from \(o_1\) to the end of \(R\).  There is no other
line with \(p\) on its left, so this is the only meaning the reading gives
to a free occurrence of \(p\) after \(o_1\).  Every free occurrence lies in
\(\alpha_1\), hence at or after \(o_1\), so each has the meaning
\(\lambda\vec z.\,s_1\).  By item~1, applied with \(n=1\), the role-(b)
reading gives each free occurrence the same meaning
\(\lambda\vec z.\,s_1\).  The occurrences bound by a statement are bound
in both readings and are not affected.

\emph{Item~3.}  Let \(n\geq2\) and suppose that a reading gives role~(a)
to \(o_i\) and \(o_j\) with \(i<j\).  By role~(a), \(o_i\) makes \(p\) a
name from \(o_i\) to the end of \(R\).  The line \(o_j\) comes after
\(o_i\), so at \(o_j\) the letter \(p\) is already a name of the region and
is not new there.  Item~2 of \cref{def:p01-w02c-p01-018-motive-reading}
requires the left side of a role-(a) line to be new in the region, and
role~(a) of \cref{def:p01-w02b-p01-013-colon-equals} requires its left side
to be a new name; so \(o_j\) cannot have role~(a).  This contradicts the
supposition, and no reading gives role~(a) to two of the lines.

\emph{Item~4.}  Let \(i<j\), let \(s_i\) and \(s_j\) be different after
every renaming of \(\vec z\), and take a reading that gives \(o_i\)
role~(a) and \(o_j\) role~(b).  By role~(a), every free occurrence of \(p\)
from \(o_i\) to the end of \(R\) means \(\lambda\vec z.\,s_i\); the
occurrence of \(p\) on the left of \(o_j\) stands after \(o_i\), so it has
this meaning.  By role~(b), the same occurrence is the letter of the rule
applied in \(\alpha_j\), and it receives the value
\(\lambda\vec z.\,s_j\) for that application.  By the comparison of
meanings in \cref{def:p01-w02c-p01-018-motive}, the two predicates are the
same meaning exactly when \(s_i\) and \(s_j\) coincide after a renaming of
\(\vec z\), which is excluded.  So that occurrence has two different
meanings.
\end{proof}
\end{annexitem}

\begin{annexitem}{annex:p01-w02c-p01-018-equality-instances}{thm:p01-w02c-p01-018-equality-instances}{Equality-eliminator uses in Harper's development}
\begin{proof}
Two facts about the calculus are used.  The eliminator of
\cref{def:p12-har-prime-regime} transports \(P\,a\) along \(a=b\) to
\(P\,b\): for \(P:A\to\mathsf{Prop}\), \(u:P\,a\) and \(c:a=b\), the term
\(E_{A,a,P}(u,b,c)\) has type \(P\,b\).  And the conversion rule of
\cref{def:p12-har-conversion} replaces a type by a convertible type
without changing the term; in particular a type \((\lambda z.\,s)\,t\)
may be replaced by its \(\beta\)-reduct \(s[t/z]\), and conversely.  The
constructor \(\mathsf{refl\_equal}\) has type \(\mathsf{eq}\ A\ a\ a\),
written \(a=a\), for every \(a:A\).

\emph{Item~1.}  The predicate \(P=\lambda z{:}A.\,(z=x)\) has type
\(A\to\mathsf{Prop}\), since \(z=x\) is \(\mathsf{eq}\ A\ z\ x\) and
\(\mathsf{eq}\ A\) has type \(A\to A\to\mathsf{Prop}\).  The type \(P\,x\)
reduces to \(x=x\), the type of \(\mathsf{refl\_equal}\ A\ x\), so by
conversion this term has type \(P\,x\).  With \(e:x=y\), the eliminator
gives \(E_{A,x,P}(\mathsf{refl\_equal}\ A\ x,y,e):P\,y\), and \(P\,y\)
reduces to \(y=x\).  By conversion the term has type \(y=x\).

\emph{Item~2.}  Here \(f\,x\) and \(f\,z\) have type \(B\) for \(z:A\), so
\(P=\lambda z{:}A.\,(f\,x=f\,z)\), with \(\mathsf{eq}\ B\), has type
\(A\to\mathsf{Prop}\).  The type \(P\,x\) reduces to \(f\,x=f\,x\), the
type of \(\mathsf{refl\_equal}\ B\ (f\,x)\); by conversion this term has
type \(P\,x\).  The eliminator with \(e:x=y\) gives a term of type
\(P\,y\), which reduces to \(f\,x=f\,y\).

\emph{Item~3.}  The predicate \(P=\lambda z{:}A.\,(x=z)\) has type
\(A\to\mathsf{Prop}\).  The type \(P\,y\) reduces to \(x=y\), the type of
\(d\), so \(d:P\,y\) by conversion.  The eliminator at the base point
\(y\), with \(e':y=w\), gives \(E_{A,y,P}(d,w,e'):P\,w\), and \(P\,w\)
reduces to \(x=w\).

\emph{Item~4.}  Read \(P\) as one name for \(\lambda z.\,(x=z)\) in all
three uses.  Then the term of item~1 is
\(E_{A,x,P}(\mathsf{refl\_equal}\ A\ x,y,e)\) with this \(P\).  The type
\(P\,x\) reduces to \(x=x\), so the seed is well typed as in item~1, and
the conclusion \(P\,y\) reduces to \(x=y\), the type of \(e\).

Read \(P\) instead as one name for \(\lambda z.\,(z=x)\) in all three uses.
The term of item~3 is \(E_{A,y,P}(d,w,e')\) with this \(P\), and the
eliminator accepts \(d\) as its seed only if \(d:P\,y\).  Since \(P\,y\)
reduces to \(y=x\), and the conversion rule passes between convertible
types in both directions, \(d:P\,y\) holds exactly when \(d\) has a type
convertible with \(y=x\).  This proves the first sentence of item~4.

For the second sentence, assume its three hypotheses.  The variable \(d\)
is declared with the type \(x=y\), and its types are exactly the types
convertible with \(x=y\).  So \(d\) has a type convertible with \(y=x\)
exactly when \(x=y\) and \(y=x\) are convertible, by transitivity and
symmetry of convertibility.  The two types \(\mathsf{eq}\ A\ x\ y\) and
\(\mathsf{eq}\ A\ y\ x\) contain no redex: \(\mathsf{eq}\) is an inductive
type and not a definition, so no \(\delta\)-step unfolds it; no
\(\lambda\)-abstraction is applied, so no \(\beta\)-step applies; and no
recursor or case analysis occurs, so no \(\iota\)-step applies.  Hence
each reduces only to itself.  Under confluence, two convertible terms have
a common reduct, which here would force
\(\mathsf{eq}\ A\ x\ y\) and \(\mathsf{eq}\ A\ y\ x\) to be the same term;
they differ in the order of the distinct variables \(x\) and \(y\).  So
they are not convertible, \(d\) does not have the type \(P\,y\), and the
term has no type.
\end{proof}
\end{annexitem}
